\documentclass[conf]{new-aiaa}
\usepackage[utf8]{inputenc}
\usepackage{graphicx}
\usepackage{amsmath}
\usepackage{siunitx}
\usepackage{longtable,tabularx}
\usepackage{booktabs}
\usepackage{float}
\usepackage[section]{placeins}
\usepackage{xcolor}
\usepackage{array}
\usepackage{xurl}
\graphicspath{{figures/}}
\emergencystretch=2.5em
\AtBeginDocument{%
  \setlength{\intextsep}{5pt plus 2pt minus 1pt}%
  \setlength{\textfloatsep}{6pt plus 2pt minus 2pt}%
  \setlength{\abovecaptionskip}{3pt}%
  \setlength{\belowcaptionskip}{1pt}%
}

\title{Layer-A Geometric-Programming Sizing of a Single Oblique-Wing Trans-Medium Vehicle: CFD-Calibrated Surrogates and a Buoyancy-Limited Payload Frontier}
\author{Yiyang I. Zhang\footnote{Independent Researcher, Winchester College.}}
\affil{Winchester College, Winchester, United Kingdom}

\begin{document}
\maketitle

\section{Methodology}

\subsection{Modelling framework and the scope of Layer~A}
The vehicle is a slender, missile-shaped hull carrying a single straight wing on one central
vertical pivot. In air the wing is deployed perpendicular to the body to generate lift; before the
vehicle enters the water the wing rotates flat against the hull, so that underwater the vehicle is a
clean body of revolution whose small residual overhang provides a hydrodynamic trim force. Sizing
such a vehicle couples three normally separate disciplines---aerodynamics, hydrodynamics, and
structures---and two operating media whose dynamic pressures differ by more than two orders of
magnitude. To keep the coupled problem tractable and, crucially, to guarantee that the optimum found
is the \emph{global} one, the design is decomposed into three layers. Layer~A, the subject of this
paper, is a geometric program (GP) that fixes every continuous sizing variable for a prescribed
wing-skew angle and a prescribed buoyancy state. Layer~B, a Bayesian optimiser that schedules the
wing-skew trajectory during the dive, sits above Layer~A and is not treated here. A computational
fluid dynamics (CFD) campaign underpins both layers by supplying the aerodynamic and hydrodynamic
coefficients that cannot be written analytically.

A geometric program is an optimisation problem in which a posynomial objective is minimised subject
to posynomial inequality constraints and monomial equality constraints. A \emph{monomial} is a
product of powers of the design variables with a positive leading coefficient; a \emph{posynomial}
is a sum of monomials. The defining property that makes the formulation attractive is that, under the
change of variables $u_i\mapsto\ln u_i$, every GP becomes convex, so a solver returns the global
optimum---not merely a local one~\cite{boyd2007tutorial}---together with the sensitivity of the objective to every
constraint. The engineering price of this guarantee is that each governing relation must be cast into
posynomial or monomial form; the methodology below is largely the discipline of paying that price
honestly.

\subsection{CFD design of experiments}
The coefficients that Layer~A cannot derive from first principles are the aerial and underwater drag
polars, the lift-curve slope and its variation with wing skew, the stowed-wing trim lift, and the
cavitation suction peak. These, and only these, define the CFD design of experiments. Two campaigns
were run in OpenFOAM~v2606 with the steady incompressible solver \texttt{simpleFoam} and the
$k$--$\omega$ SST turbulence model. The aerial campaign meshed all ten wing positions from stowed
($\Lambda=90^\circ$) to deployed ($\Lambda=0^\circ$) at $10^\circ$ intervals inside a free-air domain
of $17\times12\times12$~m, and swept the angle of attack from $-4^\circ$ to $+14^\circ$ by rotating
the inlet velocity vector rather than re-meshing, which saved roughly nine-tenths of the meshing
effort. The underwater campaign reused the same meshes with water properties, added a speed sweep of
$2$ to $20$~m/s at the deployed and stowed extremes, and treated operating depth as a
post-processing parameter, since in an incompressible solution the flow field is independent of the
absolute pressure level and depth enters only through the cavitation number. Every force coefficient
was referenced consistently to the wing planform area, and convergence was judged by coefficient
stability---the mean of the final fifty iterations was required to differ from the preceding fifty by
less than one percent---rather than by a fixed residual threshold. As an independent check on the
whole pipeline, the stowed induced-drag factor obtained in air and in water, at Reynolds numbers that
differ by a factor of about $1.5$, agreed to within $7\%$; because the drag coefficient is
dimensionless and, at matched Reynolds number, depends only on geometry, this agreement confirms that
the two campaigns are internally consistent.

\subsection{GP-compatible surrogate fitting}
Each CFD relationship was reduced to a GP-legal surrogate with the \texttt{gpfit} workflow of Hoburg~et~al.~\cite{hoburg2016data}. Where a
relationship is affine on logarithmic axes it was fitted as a single monomial by ordinary least
squares in log space, which is mathematically identical to a one-term maximum-affine fit and is
numerically robust; where a relationship carries genuine curvature, a multi-term softmax-affine
posynomial was fitted instead. The lift-curve slope varies with wing skew as
\begin{equation}
C_{L\alpha}(\Lambda) = 0.847\,(\cos\Lambda)^{1.105}\ \text{rad}^{-1},
\qquad(\text{rms}_{\log}=0.011),
\label{eq:clalpha}
\end{equation}
whose exponent of $1.105$ is within eleven percent of the value of unity predicted by the
sweep-independence principle, so the fit does more than interpolate: it confirms the analytic
expectation that only the velocity component normal to the wing generates lift. The induced-drag
factor rises steeply as the wing folds and its effective aspect ratio collapses,
\begin{equation}
k(\Lambda) = 0.631\,(\cos\Lambda)^{-1.352}\ ,
\qquad(\text{rms}_{\log}=0.013),
\label{eq:kschedule}
\end{equation}
climbing from $0.64$ at full deployment to $22.9$ when stowed---a thirty-six-fold increase that is
the direct signature of the aspect-ratio collapse in the relation $k=1/(\pi e A)$. The underwater
zero-lift drag scales with Reynolds number as $C_{D0,w}=0.185\,\mathrm{Re}^{-0.112}$, an exponent
that sits between the $-0.2$ of pure turbulent skin friction and the zero of pure form drag,
consistent with a body whose drag is a mixture of the two. The cavitation suction peak was found to
be essentially constant, as a dimensionless quantity must be: for the stowed configuration, whose
trim lift keeps it near zero lift coefficient, $-C_{p,\min}\approx1.16$, and this constant is used as
the conservative cavitation gate.

\subsection{The Layer-A geometric program}
The objective is to maximise payload, expressed as the minimisation of $m_{\text{pay}}^{-1}$. The
constraint set, transcribed from the full model of the accompanying formulation document, is written
so that every relation is a posynomial bounded by a monomial and therefore admissible in a GP. The
aerial cruise, drag build-up, and buoyancy--trim blocks are
\begin{align}
mg &\le \tfrac12\rho_a V_a^2 S\,C_{L,a}, & C_{L,a} &\le C_{L,\max}, \label{eq:aeriallift}\\
C_{D,a} &\ge \frac{C_{L,a}^2}{\pi e A} + C_{f}(1+k)\frac{S_{\text{wet}}}{S} + C_{Dp,a},
      & C_f &\ge \frac{0.074}{\mathrm{Re}^{0.2}}, \label{eq:aerialdrag}\\
(1+k) &\ge 1 + 1.5\!\left(\tfrac{d}{L}\right)^{3/2} + 7\!\left(\tfrac{d}{L}\right)^{3},
      & B &= \rho_w g\,\nabla, \label{eq:formfactor}\\
mg &\le B\,(1+\xi), & L_{\text{tr}} &\le \tfrac12\rho_w V_w^2 S\,C_{L,\text{uw}}, \label{eq:buoytrim}
\end{align}
The first relation in Eq.~\eqref{eq:aeriallift} is the steady-flight lift balance of
Hoburg and Abbeel~\cite{hoburg2014geometric}: in cruise the deployed wing must carry the full weight,
and because a larger lift coefficient always reduces the wing area the constraint is written as an
inequality that the optimiser drives to equality. It is paired with a hard ceiling
$C_{L,a}\le C_{L,\max}$ taken from the CFD stall data, without which the optimiser would demand a
physically impossible lift coefficient to shrink the wing indefinitely. Equation~\eqref{eq:aerialdrag}
is the classical drag decomposition~\cite{hoburg2014geometric,hoerner1965fluid} into an induced term,
a zero-lift viscous term, and a wing-section profile term. The induced term $C_{L,a}^2/(\pi e A)$ is
the lifting-line result~\cite{hoerner1985fluid}; the Oswald efficiency is held at the analytic prior
$e\approx0.85$ rather than the mesh-limited CFD value, for the reason set out in
Sec.~\ref{sec:justify}. The zero-lift term $C_f(1+k)S_{\text{wet}}/S$ multiplies a flat-plate
skin-friction coefficient by a form factor and a wetted-to-reference area ratio; the turbulent
power-law $C_f=0.074\,\mathrm{Re}^{-0.2}$ is used in place of the ITTC~1957 log-law because it is
monomial and its error over the relevant Reynolds range is under two
percent~\cite{hoburg2014geometric,myring1976theoretical}. Equation~\eqref{eq:formfactor} is
Hoerner's form factor for a streamlined body of revolution~\cite{hoerner1965fluid}, in which the
$1.5(d/L)^{3/2}$ term is the supervelocity increment and the $7(d/L)^3$ term the
separation-pressure component; for our fineness ratio $d/L\approx0.10$ it evaluates to
$(1+k)\approx1.054$, a five-percent viscous adder that is far below the twenty percent assumed in
early hand estimates and that we therefore adopt in place of that conservative guess. The buoyancy in
Eq.~\eqref{eq:formfactor} is Archimedes' law for the displaced hull volume $\nabla$, and the vertical
balance $mg\le B(1+\xi)$ in Eq.~\eqref{eq:buoytrim} states that weight is supported by buoyancy plus
the fraction $\xi$ carried by trim; that trim force is supplied by the stowed wing's small overhang,
$L_{\text{tr}}\le\tfrac12\rho_w V_w^2 S\,C_{L,\text{uw}}$, which is the analytic backbone of the
self-compensation thesis---the folded wing, at collapsed aspect ratio, still trims the vehicle
without a large drag penalty. The underwater drag, cavitation, energy, structure, and
feasibility-gate blocks are
\begin{align}
C_{D,w} &\ge C_{D0,w} + k_w\,C_{L,\text{uw}}^2, & D_w &\ge \tfrac12\rho_w V_w^2 S\,C_{D,w},\label{eq:hydrodrag}\\
\tfrac12\rho_w V_w^2\,(-C_{p,\min}) &\le p_{\text{atm}} + \rho_w g h - p_v,
      & E_b &\ge P_a\,t_a + P_w\,t_w, \label{eq:cavenergy}\\
t_s &\ge \frac{\rho_w g h\,d}{2\sigma_y}, & \mathrm{SF}\cdot\rho_w g h &\le 2.42\,E\!\left(\tfrac{t_s}{d}\right)^{2.5}\!\!\frac{d}{L}, \label{eq:structure}\\
S_{\text{fin}}\,l_{\text{fin}} &\ge V_{H,\min}\,S\,\bar c, & C_{l,\text{dist}} &\le C_{L\alpha,\text{fin}}\,\delta_{\max}\frac{S_{\text{fin}}}{S}\frac{b}{L}. \label{eq:gates}
\end{align}
Equation~\eqref{eq:hydrodrag} is the underwater drag polar for the stowed configuration, in the same
induced-plus-parasitic form as the aerial polar but with the induced factor $k_w$ taken directly from
the hydrodynamic campaign, because the stowed wing's aspect-ratio collapse is a real effect that no
analytic prior captures cleanly; the drag force is referenced, like every coefficient in the model,
to the wing planform area so that the CFD numbers transfer without a change of reference.
Equation~\eqref{eq:cavenergy} contains two constraints. The first is the cavitation inception
condition $-C_{p,\min}\le\sigma$ of Brennen~\cite{brennen2014cavitation}, rearranged so that the
suction pressure at the surface must not fall below the vapour pressure; Amromin and
Rozhdestvensky~\cite{amromin2022cavitation} show that at our hull Reynolds number, of order
$2\times10^{7}$, the gap between the minimum pressure coefficient and the inception number vanishes,
so the equality $\sigma_i=-C_{p,\min}$ is not merely conservative but increasingly exact, and the
constraint places a firm upper bound on underwater speed. The second is the two-leg energy balance:
for a battery-electric vehicle the fuel-burn model of a combustion aircraft is replaced by a
constant-mass balance in which the stored energy must cover the propulsive work of both the aerial
and the underwater legs~\cite{allen2000propulsion,lin2026energies}, and the battery mass follows by
dividing that energy by the cell-level specific energy. Equation~\eqref{eq:structure} sizes the
pressure-hull wall by the larger of two limits: the thin-shell hoop-stress gauge, which is monomial,
and the external-pressure buckling collapse, for which the signomial Bryant expression is replaced by
the Windenburg--Trilling closed-form monomial with a safety factor of two, following the practice
recommended for unstiffened tubes at small diameter~\cite{shinoka2024structural}.
Equation~\eqref{eq:gates} imposes the two feasibility gates. Static stability, whose exact form
involves the mixed-sign vertical-plane derivatives of Renilson~\cite{renilson2018submarine} and is
therefore signomial, is enforced through the tail-volume coefficient $V_H=S_{\text{fin}}l_{\text{fin}}
/(S\bar c)$~\cite{hoerner1985fluid}; requiring it to exceed a conventional minimum is a monomial proxy
that guarantees an adequate static margin. Roll authority requires the cruciform fins, deflecting to
their limit, to produce a rolling moment coefficient at least equal to the peak skew-induced wing
moment measured in the aerial campaign, using the low-aspect-ratio fin lift slope of
Ref.~\cite{hoerner1985fluid}. Water-entry survival is deliberately excluded from this set and applied
afterwards, for reasons given in Sec.~\ref{sec:justify}. Finally, sweeping $\xi$ keeps
Eq.~\eqref{eq:buoytrim} monomial at each solve, so the program is a pure GP throughout; freeing $\xi$
would render that equality signomial and require the sequential
solver~\cite{hoburg2014geometric,boyd2007tutorial}.

\section{Model Justification}\label{sec:justify}

\subsection{Why a geometric program, and why sweep the buoyancy state}
The choice of a geometric program is not a matter of convenience but of trust. Because the log
transform makes the problem convex, the solver certifies global optimality and reports, at no extra
cost, how hard each constraint pushes on the objective; a gradient-based optimiser on the same
non-convex problem would offer neither guarantee. The price---writing every relation as a
posynomial---is affordable here because the governing physics is already close to power-law form:
drag scales as a power of Reynolds number, buoyancy as a product of geometric dimensions, and induced
drag as the square of lift.

The buoyancy imbalance $\xi=(W-B)/B$ measures how far the vehicle's weight departs from the weight of
water it displaces. At $\xi=0$ the vehicle is neutrally buoyant and hovers with no trim force; when
$\xi>0$ it is heavy and the stowed wing must generate upward trim lift to hold depth; when $\xi<0$ it
is light and the wing trims downward. The imbalance is swept as a parameter rather than left free for
a specific reason: with $\xi$ fixed, the balance $mg=B(1+\xi)$ is monomial and the whole model
remains a pure GP; freeing $\xi$ would make that equality signomial and would require the slower
sequential (signomial-programming) solver. Sweeping $\xi$ is also the physically informative choice,
because it traces the payload-versus-buoyancy frontier directly: a heavier vehicle carries a larger
mass budget and therefore more payload, up to the point where the wing can no longer lift it in air.

\subsection{Analytic drag with CFD calibration, not a fixed CFD polar}
It would be tempting to insert the measured deployed-wing drag polar directly into the GP. That was
rejected for two reasons. First, a fixed polar is tied to one geometry; once the hull and wing become
design variables (Sec.~\ref{sec:findings}), the drag must scale with them, which only the analytic
build-up---skin friction times a form factor times a wetted-area ratio---can do. Second, and more
subtly, the deployed-wing efficiency inferred from the sweep mesh is not converged: the fitted
lift-curve slope of $0.835$~rad$^{-1}$ is about five times below the thin-airfoil expectation for the
wing's aspect ratio, and the fitted induced factor implies an Oswald efficiency of only $0.07$
against an expected $0.8$--$0.9$. Both discrepancies point to the same cause---a $57{,}000$-cell
sweep mesh under-resolving the wing's response to angle of attack---so the absolute efficiency is
better taken from the analytic prior $e\approx0.85$ than from the mesh-limited fit. What the CFD
\emph{does} deliver reliably are the dimensionless \emph{scalings} of
Eqs.~\eqref{eq:clalpha}--\eqref{eq:kschedule}, the stowed polar (whose aspect-ratio collapse is a
real effect, not a mesh artefact), the cavitation gate, and the cross-validated form factor; these
are what the model uses.

\subsection{Fixed and free geometry, the feasibility gates, and the exclusion of entry}
The model is solved in two configurations. In the first the geometry is fixed at the as-built CAD
dimensions, which verifies that the existing vehicle can meet the mission; in the second the hull
diameter and length, the wing planform, and the fin area are freed within a launch envelope, which
reveals the design the optimiser would prefer. The static-stability requirement, which in its exact
form involves stability derivatives of mixed sign and is therefore signomial, is imposed through the
standard tail-volume coefficients $V_H$ and $V_V$; requiring these to exceed conventional minima is a
GP-legal proxy that guarantees the same static-margin behaviour without leaving the convex world.
Finally, water-entry survival is kept out of the optimisation for two reasons: the peak-deceleration
coefficient is still pending its own two-dimensional impact campaign, and the impact constraint is a
pass/fail survival check that does not size the vehicle in the way the mission-energy and buoyancy
constraints do. It is therefore evaluated after the solve, as an audit of the sized design, rather
than as a driver within it.

\section{Findings}\label{sec:findings}

\subsection{CFD-fitted coefficients}
The fitted coefficients are physically self-consistent and, where an analytic prior exists, they
recover it. Figure~\ref{fig:polars} shows the drag polars for all ten skew angles: the induced-drag
factor grows smoothly and monotonically as the wing folds, with no anomalous point, and every
per-wing polar is reproduced with a coefficient of determination above $0.98$.
Figure~\ref{fig:clalpha} plots the lift-curve slope against skew and overlays the monomial fit of
Eq.~\eqref{eq:clalpha}; the near-unity exponent confirms the $\cos\Lambda$ law.
Figure~\ref{fig:kcollapse} shows the aspect-ratio collapse of the induced factor on a logarithmic
axis, and Fig.~\ref{fig:cavitation} shows why the cavitation data separate into two branches---the
stowed configuration sits at near-zero lift, while the deployed configuration spans a real lift
range---so that fitting them as a single curve, as a first attempt did, is incorrect; the stowed
branch supplies the gate value.

\begin{figure}[h]\centering
\includegraphics[width=0.86\linewidth]{aero_polars_all.png}
\caption{Aerial drag polars for all ten wing-skew angles. The induced-drag factor
$k$ (legend) rises monotonically from $0.64$ at full deployment ($\Lambda=0^\circ$) to $22.9$ when
stowed ($\Lambda=90^\circ$).}
\label{fig:polars}
\end{figure}

\begin{figure}[h]\centering
\includegraphics[width=0.72\linewidth]{sched_CLalpha.png}
\caption{Lift-curve slope versus wing skew. The fitted exponent of $1.105$ confirms the
$C_{L\alpha}\propto\cos\Lambda$ sweep-independence prior to within eleven percent.}
\label{fig:clalpha}
\end{figure}

\begin{figure}[h]\centering
\includegraphics[width=0.72\linewidth]{sched_k.png}
\caption{Induced-drag factor versus wing skew on a logarithmic axis: the thirty-six-fold rise is the
signature of the aspect-ratio collapse as the wing folds flat.}
\label{fig:kcollapse}
\end{figure}

\begin{figure}[h]\centering
\includegraphics[width=0.80\linewidth]{cavitation_split.png}
\caption{Cavitation suction peak against lift coefficient. The stowed configuration (near zero lift)
and the deployed configuration (finite lift range) form two distinct branches; the stowed branch,
$-C_{p,\min}\approx1.16$, sets the cavitation gate.}
\label{fig:cavitation}
\end{figure}

\subsection{The payload--buoyancy frontier}
Figure~\ref{fig:payload} traces payload against the buoyancy imbalance for both geometry cases. In
both, payload rises with $\xi$, because a heavier vehicle carries a larger mass budget; in both, the
frontier terminates where the aerial lift coefficient reaches its stall ceiling and the wing can no
longer support the vehicle in air. The fixed-geometry vehicle carries between $2.4$ and $3.5$~kg of
payload across the feasible range of $\xi$, which confirms that the as-built design meets a several-
kilogram payload goal. Freeing the geometry roughly doubles the payload, to between $5.9$ and
$8.4$~kg, for a reason that Sec.~\ref{sec:optgeom} makes precise: the optimiser enlarges the hull to
the edge of the launch envelope and thereby buys buoyancy.

\begin{figure}[h]\centering
\includegraphics[width=0.72\linewidth]{payload_xi.png}
\caption{Payload versus buoyancy imbalance $\xi$ for the fixed CAD geometry and the
envelope-constrained free geometry. Both curves end where the aerial lift coefficient reaches the
$0.60$ stall ceiling.}
\label{fig:payload}
\end{figure}

\subsection{The optimised geometry and what constrains it}\label{sec:optgeom}
Table~\ref{tab:soln} lists the free-geometry solution across the imbalance sweep. The pattern is
uniform and revealing: at every value of $\xi$ the hull diameter and length pin to their
envelope caps of $0.12$~m and $1.30$~m, the wing span pins to the geometric limit of one hull length
minus $0.30$~m, and the aerial lift coefficient pins to its $0.60$ stall ceiling. In other words the
optimiser drives the vehicle to three limits at once---how large the launch platform allows it to be,
how wide the fuselage allows the wing to span, and how hard the wing can work in air---and the
payload it returns is whatever remains after the mission energy and the pressure hull have been paid
for. The cruciform fins, sized by the static-stability gate, come out at about $86$~cm$^2$ of total
area on a $0.585$~m moment arm, adding roughly three-quarters of a kilogram; the roll-authority gate
is satisfied with margin by the fins that stability already demands. Neither gate drives the payload,
which is exactly the behaviour a feasibility gate should show.

\begin{table}[h]\centering
\caption{Free-geometry Layer-A solution across the buoyancy sweep. The row at $\xi=0.15$ is the
reference design point.}\label{tab:soln}
\small
\begin{tabular}{@{}rrrrrrrr@{}}\toprule
$\xi$ & payload & mass & $d$ & $L$ & span & $S$ & $S_{\text{fin}}$ \\
      & (kg) & (kg) & (mm) & (mm) & (m) & (m$^2$) & (cm$^2$) \\\midrule
$0.00$ & $5.86$ & $9.12$ & $120$ & $1300$ & $1.00$ & $0.097$ & $64.8$ \\
$0.05$ & $6.28$ & $9.57$ & $120$ & $1300$ & $1.00$ & $0.102$ & $71.4$ \\
$0.10$ & $6.71$ & $10.03$ & $120$ & $1300$ & $1.00$ & $0.107$ & $78.4$ \\
$0.15$ & $7.14$ & $10.48$ & $120$ & $1300$ & $1.00$ & $0.112$ & $85.7$ \\
$0.20$ & $7.56$ & $10.94$ & $120$ & $1300$ & $1.00$ & $0.117$ & $93.3$ \\
$0.25$ & $7.99$ & $11.39$ & $120$ & $1300$ & $1.00$ & $0.122$ & $101.2$ \\
$0.30$ & $8.41$ & $11.85$ & $120$ & $1300$ & $1.00$ & $0.127$ & $109.5$ \\\bottomrule
\end{tabular}
\end{table}

\subsection{Sensitivities and design levers}
The dual solution ranks the levers that matter. The strongest is the displaced volume, with a
sensitivity of $-1.43$: enlarging the hull is the most direct way to increase payload, which is why
the optimiser pins the hull to its bounds. Battery specific energy follows at $-0.49$, then aerial
efficiency and range at about $\pm0.30$, and required underwater speed at $+0.38$. The stowed
induced-drag factor registers a sensitivity of only $+0.001$---negligible---because the stowed wing
operates at almost zero lift and its trim drag is correspondingly tiny. This last number is worth
dwelling on, because it is the quantitative form of the vehicle's design thesis: the folded wing
costs almost nothing in drag, so the penalty for keeping the wing rather than discarding it before
water entry is small.

\subsection{Verification of the surrogates and the sized model}
Each surrogate and the assembled program were checked by five independent means before either was
trusted. The first is the air--water cross-validation already noted: the stowed induced-drag factor
agrees to seven percent between the two campaigns, a comparison that no single campaign could make.
The second is agreement with analytic priors: the lift-slope exponent of $1.105$ against the value of
unity in Eq.~\eqref{eq:clalpha}, the Reynolds exponent of $-0.112$ bracketed by the skin-friction and
form-drag limits, and the near-zero exponent of the cavitation peak all fall where first principles
say they should. The third is reconstruction quality: every per-wing polar is recovered with a
coefficient of determination above $0.98$ and every retained monomial fit with a root-mean-square
logarithmic error below $0.02$. The fourth is GP-legality, verified term by term: each quantity that
enters a constraint has a positive-coefficient form, the one exception---the cambered deployed wing,
whose zero-lift intercept is negative---being replaced by a positive softmax-affine surrogate. The
fifth is physical monotonicity: lift rises with angle of attack, drag rises as the wing deploys, and
the suction peak is speed-independent, at every point sampled. The assembled program was audited in
the same spirit: its dual sensitivities are all of the expected sign and magnitude, and its binding
constraints---mass closure and buoyancy---are exactly those a near-neutral trans-medium vehicle
should be limited by, which is itself evidence that the model is wired correctly.

\section{Layer-A Conclusions}

Three conclusions follow from the Layer-A study. First, the vehicle is buoyancy- and
envelope-limited rather than aerodynamically or structurally limited: the binding constraints are the
hull-size caps and the aerial stall ceiling, and the dominant sensitivity is to displaced volume, so
payload is bought almost directly with the size the launch platform will allow. The as-built CAD
vehicle carries roughly three kilograms of payload; an envelope-optimised vehicle of the same family
carries about seven. Second, the design thesis survives contact with the numbers: the stowed wing's
trim drag is negligible in the sizing, so retaining a steerable wing through water entry, rather than
jettisoning it, imposes almost no payload penalty, and the lift-curve slope's clean $\cos\Lambda$
scaling confirms that the oblique wing behaves as classical sweep theory predicts. Third, the model
is honest about where it is weakest: the absolute efficiency of the deployed wing is not yet
grid-converged and is therefore carried through an analytic prior rather than the mesh-limited fit,
the drag polar over the measured lift range is too narrow to support a full softmax-affine surrogate,
and the water-entry constant awaits its own campaign. These are the three inputs a refined study
should target---a grid-convergence mesh for the deployed wing, a wider angle-of-attack sweep or a
section-level profile-drag polar, and the two-dimensional impact campaign---after which the same
Layer-A program, unchanged in structure, will return sharper numbers. The value of the present model
is that it already tells the designer where to look: at the size of the launch envelope, the energy
density of the battery, and the efficiency of the deployed wing, in that order.

% \bibliographystyle is set by new-aiaa.cls
\bibliography{uauv}

\end{document}
